\section{RNN 的内部机制}

\subsection{RNN 的直觉理解}

讲者通过一段 Python 伪代码来建立对 RNN 的直觉理解：

\begin{figure}[H]
\centering
\includegraphics[width=0.85\textwidth]{slides-images/slide-20.jpg}
\caption{RNN 伪代码直觉：遍历句子中的每个词，更新隐藏状态并生成预测\protect\footnotemark}
\end{figure}
\footnotetext{Slides 第20页。}

\begin{lstlisting}[caption={RNN 伪代码示例}]
my_rnn = RNN()
hidden_state = [0, 0, 0, 0]

sentence = ["I", "love", "recurrent", "neural"]

for word in sentence:
    prediction, hidden_state = my_rnn(word, hidden_state)

next_word_prediction = prediction
# >>> "networks!"
\end{lstlisting}

这段伪代码清楚展示了 RNN 的工作流程：
\begin{enumerate}
    \item 初始化隐藏状态为零向量
    \item 对序列中的每个元素，调用 RNN 处理当前输入和上一步的隐藏状态
    \item 每次调用返回一个预测和更新后的隐藏状态
    \item 序列处理完毕后，最终的预测即为输出
\end{enumerate}

\subsection{RNN 状态更新的数学公式}

\begin{figure}[H]
\centering
\includegraphics[width=0.75\textwidth]{slides-images/slide-22.jpg}
\caption{RNN 状态更新与输出机制\protect\footnotemark}
\end{figure}
\footnotetext{Slides 第22页。}

RNN 的前向传播包含两个关键步骤：

\textbf{Step 1: 更新隐藏状态}

$$h_t = \tanh(W_{xh} \cdot x_t + W_{hh} \cdot h_{t-1} + b_h)$$

\begin{itemize}
    \item $W_{xh}$：将输入转化到隐藏空间的权重矩阵
    \item $W_{hh}$：隐藏状态自身更新的权重矩阵
    \item $b_h$：偏置项
    \item $\tanh$：非线性激活函数
\end{itemize}

\textbf{Step 2: 生成输出}

$$\hat{y}_t = W_{hy} \cdot h_t + b_y$$

\begin{itemize}
    \item $W_{hy}$：将隐藏状态转化为输出的权重矩阵
    \item $b_y$：输出偏置项
\end{itemize}

\begin{importantbox}{权重共享：RNN 的关键特性}
在 RNN 中，权重矩阵 $W_{xh}$、$W_{hh}$ 和 $W_{hy}$ 在\textbf{所有时间步上共享}。这意味着无论序列多长，参数数量不变。这不仅大大减少了模型参数，还使 RNN 能够处理任意长度的序列。
\end{importantbox}

\subsection{RNN 的计算图展开}

\begin{figure}[H]
\centering
\includegraphics[width=0.85\textwidth]{slides-images/slide-25.jpg}
\caption{RNN 计算图在时间轴上的展开\protect\footnotemark}
\end{figure}
\footnotetext{Slides 第25页。}

将 RNN 在时间轴上展开后，可以清楚地看到：
\begin{itemize}
    \item 每个时间步接收输入 $x_t$，通过 $W_{xh}$ 参与隐藏状态计算
    \item 隐藏状态 $h_t$ 通过 $W_{hh}$ 传递到下一个时间步
    \item 每个时间步通过 $W_{hy}$ 产生输出预测 $\hat{y}_t$
\end{itemize}

\subsection{RNN 的代码实现}

\begin{lstlisting}[caption={基于 TensorFlow 的 RNN 层实现}]
class MyRNNCell(tf.keras.layers.Layer):
    def __init__(self, rnn_units):
        super(MyRNNCell, self).__init__()
        self.W_xh = self.add_weight([rnn_units, rnn_units])
        self.W_hh = self.add_weight([rnn_units, rnn_units])
        self.W_hy = self.add_weight([rnn_units, num_outputs])

    def call(self, x, h):
        # Update hidden state
        h = tf.math.tanh(
            tf.linalg.matmul(x, self.W_xh) +
            tf.linalg.matmul(h, self.W_hh)
        )
        # Compute output
        output = tf.linalg.matmul(h, self.W_hy)
        return output, h
\end{lstlisting}

在 TensorFlow/Keras 等框架中，也可以直接使用封装好的 RNN 层：

\begin{lstlisting}[caption={使用 Keras 封装的 SimpleRNN 层}]
tf.keras.layers.SimpleRNN(rnn_units)
\end{lstlisting}

\subsection{本章小结}

RNN 的核心操作可以归纳为三步：(1) 接收当前输入；(2) 结合上一步隐藏状态更新当前隐藏状态；(3) 基于隐藏状态产生输出。权重在所有时间步共享是 RNN 能够处理变长序列的关键。

\newpage
